\documentclass[11pt]{article}
\usepackage{cmap}
\usepackage[T1]{fontenc}
\usepackage[utf8]{inputenc}
\usepackage{lmodern}
\usepackage{amsmath,amssymb}
\usepackage{microtype}
\usepackage[margin=1in]{geometry}
\usepackage[hidelinks]{hyperref}
\hypersetup{
  pdftitle={Solving SAT in polynomial time?},
  pdfauthor={Brian Tenneson},
  pdfsubject={SAT and Groebner bases},
  pdfkeywords={P vs NP, polynomial time SAT solver, Boolean algebra}
}

\title{Solving SAT in polynomial time?}
\author{Brian Tenneson}
\date{March 22, 2025}

\begin{document}
\maketitle

\begin{abstract}
Our goal is to shed some light on the $\mathcal{P}$ vs $\mathcal{NP}$ problem by looking at a SAT solver and investigate whether it runs in polynomial time. Given the $\mathcal{NP}$-completeness of SAT, if it does, this would imply $\mathcal{P}=\mathcal{NP}$.
\end{abstract}

\noindent\textbf{Keywords:} $\mathcal{P}$ vs $\mathcal{NP}$, polynomial time SAT solver, Boolean algebra

\medskip
\noindent We would like to acknowledge all of the many people who contributed to this work.

\section{A SAT solver which might run in polynomial time}

\subsection{Input Boolean wff $\varphi$}
Either $\varphi=p_i$ for a propositional variable $p_i$ or $\varphi=\neg\varphi_1$ for some Boolean wff $\varphi_1$, or $\varphi=\varphi_1\wedge\varphi_2$ for Boolean wffs $\varphi_1$ and $\varphi_2$, or $\varphi=\varphi_1\vee\varphi_2$ for Boolean wffs $\varphi_1$ and $\varphi_2$.

\subsection{Compute a polynomial associated with $\varphi$ called $\tau(\varphi)$ as follows:}
\begin{align*}
\tau(p_i) &:= x_i,\\[0.65em]
\tau(\neg\varphi_1) &:= 1+\tau(\varphi_1),\\[0.65em]
\tau(\varphi_1\wedge\varphi_2) &:= \tau(\varphi_1)\tau(\varphi_2),\\[0.65em]
\tau(\varphi_1\vee\varphi_2) &:= 1+\bigl(1+\tau(\varphi_1)\bigr)\bigl(1+\tau(\varphi_2)\bigr).
\end{align*}
Calculations are done in $\mathbb{F}_2$.

\subsection{Let $J=\langle \tau+1,x_1^2+x_1,\ldots,x_n^2+x_n\rangle$ be an ideal in $\mathbb{F}_2[x_1,\ldots,x_n]$}
Note that
\[
J=\langle\tau+1\rangle+\langle x_1^2-x_1,\ldots,x_n^2-x_n\rangle.
\]

\newpage
\subsection{Find a Gr\"obner basis $G$ for $J$ which might be done in polynomial time using the FGLM algorithm or a different algorithm}

\subsection{Check whether $\overline{1}^{\,G}=0$}

\subsubsection{If $\overline{1}^{\,G}=0$, then $\varphi$ is not satisfiable}

\subsubsection{If $\overline{1}^{\,G}\neq 0$, then $\varphi$ is satisfiable}

\end{document}
